\documentclass[10pt,a4paper]{amsart}
\usepackage[margin=19mm]{geometry}
\usepackage{amsmath,amssymb,mathtools}
\usepackage[scr=rsfs,cal=euler]{mathalfa}
\usepackage{xfrac}
\usepackage[unicode,hidelinks,bookmarks=false]{hyperref}
\newcommand{\ie}{i.e.\ }
\newcommand{\cf}{cf.\ }
\newcommand{\resp}{respectively\ }
\newcommand{\Subsection}{\subsection}
\providecommand{\nobreakdash}{\nobreak\hbox{-}}
\setlength{\emergencystretch}{2em}
\multlinegap0pt
\usepackage{algorithm}\usepackage{algpseudocode}
\usepackage{amssymb}
\usepackage{mathtools}
\usepackage{dsfont}
\usepackage{xcolor}
\usepackage{enumerate}
\usepackage{bm}
\usepackage[shortlabels]{enumitem}
\usepackage{float}
\usepackage{csquotes}
\newtheorem{proposition}{Proposition}
\usepackage{cleveref}
\crefname{proposition}{Proposition}{Propositions}
\newcommand{\bu}{{\bm{u}}}
\newcommand{\bv}{{\bm{v}}}
\newcommand{\bz}{{\bm{z}}}
\newcommand{\norm}[1]{\left\lVert #1 \right\rVert}
\title{Iterated graph Laplacian for image restoration problems}
\author{Stefano Aleotti, Davide Bianchi, Florian Bossmann, Marco Donatelli, Pietro Maurino}
\date{Source: arXiv:2607.17313v1 (CC BY 4.0)}
\begin{document}
\maketitle

\noindent\emph{Source and licence.} Iterated graph Laplacian for image restoration problems. Version arXiv:2607.17313v1 (CC BY 4.0), distributed by arXiv under the Creative Commons Attribution 4.0 International licence, http://creativecommons.org/licenses/by/4.0/, as recorded in the arXiv licence metadata for this exact version and shown on the abstract page https://arxiv.org/abs/2607.17313v1. Full licence notice: \texttt{licenses/arxiv-2607.17313-CC-BY-4.0.txt}.\par\smallskip

\noindent\textit{Editorial scope.} Counted unit: the article's proposition prop:dim-free-deblurring (source0.tex lines 2132--2138). The article's complete original proof of it is delivered with it (source0.tex lines 2140--2201, one \texttt{proof} environment of 62 lines). No mathematics is written, repaired or re-derived: the statement and the proof are the article's own lines, in the article's own notation, and no retained line is substituted or re-flowed.  Retained uncounted as the source-local context the proof needs: lines 2126--2129.  Nothing was omitted from any retained span, so the package carries no omission record.  No retained line contains a citation, so the wrapper carries no bibliography and no bibliography entry.  No retained line contains a figure environment, an image-inclusion command or drawing code, so no image is delivered and the package declares no asset dependency.  Editorial formatting only, in addition: an AMS \texttt{amsart} preamble that restates the article's own declarations and macros used by the retained lines, namely the environments \texttt{proposition} on separate counters, together with the standard packages those lines need.  The retained environments print in this document as Proposition 1 in order of appearance, whereas the article prints the same items with the numbering of its own full document.  The complete native source of the article as distributed in the arXiv e-print for this exact version is bundled unchanged as \texttt{sources/205555/source0.tex}.
\begin{equation}\label{eq:graph-form-comparison}
\langle \Delta_\bz\bu,\bu\rangle
\ge \omega_*\langle L_N\bu,\bu\rangle,
\end{equation}

\begin{proposition}\label{prop:dim-free-deblurring}
Under the assumptions above, there exists a constant $C_J>0$, independent of $N$ and of $\bz$, such that
\[
\norm{\bu}_2\le C_J\bigl(\norm{K_N\bu}_2+\norm{\Delta_\bz\bu}_2\bigr)
\qquad\text{for every }\bu\in X_N.
\]
\end{proposition}

\begin{proof}
Since $\widehat h$ is continuous and $\widehat h(0)\ne0$, there exist $\tau>0$ and $\kappa>0$, independent of $N$, such that
\[
|\widehat h(\theta)|\ge \kappa
\qquad\text{whenever }\lambda(\theta)\le\tau.
\]
Let $P_\tau$ be the Fourier projection onto the modes with $\lambda(\theta)\le\tau$, and write
\[
\bu_\ell=P_\tau\bu,
\qquad
\bu_h=(I-P_\tau)\bu.
\]
On the low-frequency space, $K_N$ is uniformly bounded below:
\[
\norm{\bu_\ell}_2\le \kappa^{-1}\norm{K_N\bu_\ell}_2.
\]
By Young's inequality, $\norm{K_N\bv}_2\le M_K\norm{\bv}_2$ with $M_K=\norm{h}_{\ell^1}$, so
\begin{equation}\label{eq:low-bound-blur}
\norm{\bu_\ell}_2
\le \kappa^{-1}\norm{K_N\bu}_2+\kappa^{-1}M_K\norm{\bu_h}_2.
\end{equation}
For the high-frequency part, the definition of $P_\tau$ gives
\[
\tau\norm{\bu_h}_2^2\le \langle L_N\bu,\bu\rangle.
\]
Using \Cref{eq:graph-form-comparison} and Cauchy--Schwarz,
\begin{equation}\label{eq:high-bound-blur}
\norm{\bu_h}_2^2
\le (\tau\omega_*)^{-1}\langle\Delta_\bz\bu,\bu\rangle
\le (\tau\omega_*)^{-1}\norm{\Delta_\bz\bu}_2\norm{\bu}_2.
\end{equation}
Combining \Cref{eq:low-bound-blur} with $\norm{\bu}_2\le\norm{\bu_\ell}_2+\norm{\bu_h}_2$ gives
\[
\norm{\bu}_2
\le \kappa^{-1}\norm{K_N\bu}_2
+\bigl(1+\kappa^{-1}M_K\bigr)\norm{\bu_h}_2.
\]
Insert \Cref{eq:high-bound-blur}. With
\[
a:=\norm{\bu}_2,
\qquad b:=\norm{K_N\bu}_2,
\qquad c:=\norm{\Delta_\bz\bu}_2,
\]
we obtain
\[
a\le \kappa^{-1}b+
\bigl(1+\kappa^{-1}M_K\bigr)(\tau\omega_*)^{-1/2}(ca)^{1/2}.
\]
Set
\[
A:=\bigl(1+\kappa^{-1}M_K\bigr)(\tau\omega_*)^{-1/2}.
\]
Young's inequality gives $A(ca)^{1/2}\le \frac12 a+\frac12 A^2c$. Therefore
\[
a\le 2\kappa^{-1}b+A^2c.
\]
Equivalently,
\[
\norm{\bu}_2\le C_J\bigl(\norm{K_N\bu}_2+\norm{\Delta_\bz\bu}_2\bigr),
\]
with $C_J=\max\{2\kappa^{-1},A^2\}$, which is independent of $N$ and $\bz$.
\end{proof}

\end{document}
